\documentclass[11pt,a4paper,sans]{moderncv}

% ModernCV themes
\moderncvstyle{banking}
\moderncvcolor{blue}

% Character encoding
\usepackage[utf8]{inputenc}

% Adjust page margins
\usepackage[scale=0.82]{geometry}
\setlength{\hintscolumnwidth}{3cm}

% Personal data
\name{Md Rony}{Golder}
\title{PhD}
\address{Perth, Western Australia, Australia}{}{}
\phone[mobile]{Available upon request}
\email{mdrony.golder@uwa.edu.au}
\email{ronygolderku@gmail.com}
\homepage{ronygolderku.netlify.app}
\social[linkedin]{ronygolder}
\social[github]{ronygolderku}
\extrainfo{ORCID: 0000-0002-6863-139X}

\begin{document}

\makecvtitle

\section{Professional Summary}
Research Fellow at The University of Western Australia specializing in ocean color remote sensing, biological oceanography, and marine ecosystem modeling. PhD in Spatial Science with 27 research outputs including 18 peer-reviewed articles (h-index: 7, 95 citations - Scopus). Core expertise: water quality assessment, TUFLOW/AED model validation, satellite data analysis, GIS, and marine biogeochemical processes. Previously: Consultant at BMT, Research Officer at UWA, Data Analyst at WAMSI/SEAF, Casual Academic at Curtin University, Visiting Fellow at Université de Brest, and Research Intern at CSIRO.

\section{Education}
\cventry{2022--2026}{Doctor of Philosophy in Spatial Science}{Curtin University}{Bentley, WA, Australia}{}{
  \textbf{Degree Conferred:} 2026\\
  \textbf{Thesis Title:} Long term changes of phytoplankton in the Southern Ocean\\
  \textbf{Supervisors:} Dr. David Antoine (Principal), Dr. Peter Fearns, Dr. Alex Sen Gupta, Dr. Pete Strutton\\
  \textbf{Academic Status:} PhD degree conferred 2026
}

\cventry{2019--2020}{Master of Science in Coastal and Marine Science}{Khulna University}{Khulna, Bangladesh}{\textbf{CGPA: 3.83 (Distinction)}}{
  \textbf{Merit Position:} 1st Class 1st\\
  \textbf{Thesis:} Seasonal and annual distribution of particulate organic carbon (POC) in the Bay of Bengal using satellite imagery\\
  \textbf{Thesis available online:} \url{https://ronygolderku.github.io/ms_thesis_rony/}
}

\cventry{2015--2018}{Bachelor of Science in Fisheries (Honours)}{Khulna University}{Khulna, Bangladesh}{\textbf{CGPA: 3.80 (Distinction)}}{
  \textbf{Merit Position:} 1st Class 2nd\\
  \textbf{Thesis:} Effects of different experimental diets and salinity on growth performance of \textit{Penaeus monodon}\\
  \textbf{Thesis available online:} \url{https://ronygolderku.github.io/BSc_thesis_rony/}
}

\section{Professional Experience}

\cventry{Jun 2026--Present}{Research Fellow}{The University of Western Australia}{Perth, WA, Australia}{Full-time · On-site}{
  \begin{itemize}
    \item Working on water quality assessment of Darwin Harbour using satellite and in-situ data
    \item TUFLOW model validation, data analysis, and visualisation for coastal ecosystem studies
    \item Support research projects related to coastal and marine ecosystems
    \item Collaborate with Centre for Water and Spatial Science on marine modeling projects
  \end{itemize}
  \textbf{Skills:} Water quality assessment, TUFLOW, Model validation
}

\cventry{Dec 2025--Jun 2026}{Research Officer}{The University of Western Australia}{Perth, WA, Australia}{Part-time · Hybrid}{
  \begin{itemize}
    \item Coordinated with the Centre for Water and Spatial Science on AED model quality assurance
    \item Validated model outputs using field observations and satellite data
    \item Analysed chlorophyll-a, phytoplankton, TSS, and CDOM to assess ecosystem dynamics
    \item Developed time-series and profile analyses for productivity and biogeochemical interpretation
    \item Contributed to Cockburn Sound Integrated Ecosystem Model development
  \end{itemize}
  \textbf{Skills:} AED (Aquatic EcoDynamics), Biogeochemistry, Time-series analysis
}

\cventry{Jan 2026--May 2026}{Consultant}{BMT}{Perth, WA, Australia}{Part-time · On-site}{
  \begin{itemize}
    \item Evaluated BMT-developed numerical model performance by validating against satellite and in situ observations
    \item Analysed hydrodynamics, water quality, and biogeochemical processes for investigative and applied studies
    \item Post-processed and analysed TUFLOW FV model outputs using Python
    \item Supported data handling, visualisation, and interpretation of model results
    \item Contributed to client reporting and technical documentation
  \end{itemize}
  \textbf{Skills:} Numerical modelling, TUFLOW FV, Python, Water quality analysis
}

\cventry{Mar 2024--Jan 2026}{Research Assistant / Data Analyst}{WAMSI/SEAF}{Perth, WA, Australia}{Part-time · Remote}{
  \begin{itemize}
    \item Developed automated data pipelines for environmental monitoring programs
    \item Processed and quality-controlled oceanographic datasets from multiple sources
    \item Created interactive dashboards for marine data visualization using Python and R
    \item Collaborated with stakeholders on data management strategies
    \item Implemented version control and documentation standards for reproducible research
  \end{itemize}
  \textbf{Skills:} Data pipelines, Quality control, Dashboard development, Stakeholder engagement
}

\cventry{May 2025--Jun 2025}{Visiting Fellow}{Université de Brest}{Brest, France}{International Doctoral Mobility Scholarship}{
  \begin{itemize}
    \item Applied deep learning techniques to phytoplankton reconstruction from ocean color data
    \item Collaborated with LOPS (Laboratoire d'Océanographie Physique et Spatiale) researchers
    \item Developed neural network models for satellite data gap-filling
    \item Presented research findings at LOPS seminar series
  \end{itemize}
  \textbf{Funding:} €2,750 (Erasmus+) + A\$3,200 (Curtin University)\\
  \textbf{Skills:} Deep learning, Neural networks, Ocean color remote sensing
}

\cventry{Aug 2024--Nov 2024}{Research Intern}{CSIRO (Oceans and Atmosphere)}{Hobart, TAS, Australia}{}{
  \begin{itemize}
    \item Integrated biogeochemical observational data into Bluelink ReANalysis (BRAN) models
    \item Analysed satellite-derived chlorophyll-a and in-situ measurements
    \item Validated model outputs against Argo float and ship-based observations
    \item Contributed to CSIRO's operational oceanography program
  \end{itemize}
  \textbf{Skills:} BRAN models, Model-data integration, Argo float data processing
}

\cventry{Feb 2026--May 2026}{Casual Academic}{Curtin University}{Perth, WA, Australia}{Teaching}{
  \begin{itemize}
    \item Taught SPAT2008 – Environmental Applications of Geographic Information Science and Remote Sensing
    \item Delivered practical sessions on ArcGIS Pro, Python for GIS, and satellite data analysis
    \item Supervised student projects on environmental remote sensing applications
    \item Unit Coordinator: Dr. Todd Robinson
  \end{itemize}
  \textbf{Skills:} Teaching, GIS instruction, Remote sensing education
}

\section{Research Metrics (July 2026)}
\cvitem{Total Outputs}{27 research outputs (18 peer-reviewed articles, 6 review articles, 2 theses, 1 press article)}
\cvitem{H-Index}{7 (Scopus)}
\cvitem{Citations}{95 (Scopus)}
\cvitem{Research Areas}{Ocean Color Remote Sensing • Biological Oceanography • GIS • Marine Ecology}
\cvitem{Profiles}{Google Scholar: \url{https://scholar.google.com/citations?user=yXg7upIAAAAJ}}
\cvitem{}{Scopus Author ID: 57226775186}
\cvitem{}{ORCID: 0000-0002-6863-139X}

\section{Recent Achievements}

\subsection{Major Research Grants (2026)}
\cvitem{ACEAS Impact Project Grant}{\textbf{A\$56,112} (March--June 2026)\\
  \textit{Project:} Open Access digital resource dedicated to Earth Observation (remote sensing) and in-situ data\\
  \textit{Role:} Content Lead \& Principal Investigator\\
  \textit{Funder:} Australian Centre for Excellence in Antarctic Science (ACEAS)\\
  \textit{Objective:} Develop comprehensive online resource for Antarctic researchers
}

\section{Honours \& Awards (Total Funding: A\$250,000+)}

\subsection{Major Research Grants}
\cvitem{2022--2026}{Curtin International Postgraduate Research Scholarship (CIPRS)\\
  \textbf{A\$37,500/year × 3.5 years = A\$131,250}\\
  Curtin University, Australia}

\cvitem{2025}{International Doctoral Mobility Scholarship\\
  \textbf{€2,750 (Erasmus+) + A\$3,200 (Curtin) = A\$7,400}\\
  Université de Brest, France \& Curtin University, Australia}

\cvitem{2024}{AGU Student Travel Grant\\
  \textbf{US\$1,000}\\
  American Geophysical Union\\
  To attend Ocean Sciences Meeting 2026, Glasgow, Scotland}

\cvitem{2023}{ACEAS Travel Grant\\
  \textbf{A\$2,570}\\
  Australian Centre for Excellence in Antarctic Science\\
  For conference attendance in Hobart, Australia}

\subsection{International Fellowships}
\cvitem{2017}{DAAD Scholarship (Deutscher Akademischer Austauschdienst)\\
  Leibniz-Institut für Ostseeforschung Warnemünde (IOW), Germany\\
  Research visit to study Baltic Sea oceanography}

\cvitem{2012--2020}{National Science and Technology (NST) Fellowship\\
  Ministry of Science and Technology, Government of Bangladesh\\
  Awarded for academic excellence in science}

\subsection{Competitive Awards}
\cvitem{2020}{Best Presenter Award\\
  7th Bangladesh Fisheries Conference, Bangladesh\\
  For presentation on ocean color remote sensing applications}

\cvitem{2019}{OCRSD Training Scholarship\\
  Indian National Centre for Ocean Information Services (INCOIS), India\\
  Fully funded training on Ocean Colour Remote Sensing – Data, Processing and Applications}

\subsection{Academic Scholarships}
\cvitem{2019--2020}{Master's Merit Scholarship\\
  Khulna University, Bangladesh\\
  Awarded for 1st Class 1st position in MSc}

\cvitem{2015--2018}{Undergraduate Merit Scholarship\\
  Khulna University, Bangladesh\\
  Awarded for academic excellence}

\cvitem{2018}{Bangabandhu Science and Technology Fellowship\\
  Ministry of Science and Technology, Government of Bangladesh}

\cvitem{2012}{General Grade Scholarship\\
  Bangladesh Education Board\\
  For Secondary School Certificate (SSC) results}

\section{Publications}

\subsection{Manuscripts Under Review (2)}
\begin{enumerate}
  \item \textbf{Golder, M. R.}, \& Antoine, D. (2026). Antarctic Sea-ice Extent Anomalies Impact Interannual Variability of Phytoplankton Chlorophyll-a in Southern Ocean Frontal Zones. \textit{Geophysical Research Letters}. (Under Review)

  \item Shammi, A. T., et al. (2026). Assessing Environmental Services and Recreational Value of Jashore Pouro Park In Jashore District of Bangladesh. \textit{Heliyon}. (Under Review)
\end{enumerate}

\subsection{Peer-Reviewed Publications}

\textbf{2026 (1 publication)}
\begin{enumerate}
  \item Thornber, K., Haque, M. M., Rouf, M. A., Fardoshi, T., Hasan, N. A., \textbf{Golder, M. R.}, Bashar, A., Debnath, P. P., Rahman, M. M., Kawahigashi, D., Chadag, V. M., \& Tyler, C. R. (2026). Improving sustainable water usage in extensive shrimp aquaculture in Bangladesh. \textit{Aquaculture}, 624, 744134. \url{https://doi.org/10.1016/j.aquaculture.2026.744134}
\end{enumerate}

\textbf{2025 (3 publications)}
\begin{enumerate}
  \item \textbf{Golder, M. R.}, \& Antoine, D. (2025). Physical drivers of long-term chlorophyll-a variability in the Southern Ocean. \textit{Elementa: Science of the Anthropocene}, 13(1), 00077. \url{https://doi.org/10.1525/elementa.2024.00077}

  \item Hipsey, M., Huang, P., Busch, B., Zhai, S., \textbf{Golder, M. R.}, \& Paraska, D. (2025). The Cockburn Sound Integrated Ecosystem Model. \textit{The University of Western Australia}. \url{https://doi.org/10.26182/tp2x-t920}

  \item Islam, S. S., Nowshin, R., Bir, J., Sarmin, S. S., Islam, M. R., Sultana, A., Ferdose, J., Sultana, A., Zohora, F. T., Khanom, M., \textbf{Golder, M. R.}, Zobayer, M. F. A., Das, K. K., Khatun, M. A., Sadia, S., \& Huq, K. A. (2025). Impact of commercial probiotics on growth, prophenoloxidase, superoxide dismutase, and digestive enzyme activity in \textit{Macrobrachium rosenbergii} cultured at different stocking densities. \textit{Aquaculture Research}, 1, 4397755. \url{https://doi.org/10.1155/are/4397755}
\end{enumerate}

\textbf{2023 (4 publications)}
\begin{enumerate}
  \item Hosen, M., Alam, M., Chakraborty, T., \& \textbf{Golder, M. R.} (2023). Monitoring spatiotemporal and seasonal variation of agricultural drought in Bangladesh using MODIS-derived vegetation health index. \textit{Journal of Earth System Science}, 132, 188. \url{https://doi.org/10.1007/s12040-023-02200-3}

  \item Shammi, A. T., Hassan, N., \textbf{Golder, M. R.}, Molla, H., \& Islam, S. S. (2023). Health status assessment of people adjacent to temporary waste disposal sites in Khulna city, Bangladesh. \textit{Heliyon}, 9(9). \url{https://doi.org/10.1016/j.heliyon.2023.e19810}

  \item Fardoshi, T., \textbf{Golder, M. R.}, Rouf, M. A., \& Masud-Ul-Alam, M. (2023). Seasonal and spatial variability of chlorophyll-a in response to ENSO and ocean current in the maritime boundary of Bangladesh. \textit{Journal of Marine Sciences}, 3(12). \url{https://doi.org/10.1155/2023/2843608}

  \item Munia, T. T., Islam, H. M. R., \textbf{Golder, M. R.}, Hossen, M. B., Banu, G. R., \& Islam, S. S. (2023). Response of \textit{Macrobrachium rosenbergii} juveniles against a virulent isolate of \textit{Vibrio} sp. from a diseased \textit{Penaeus monodon}. \textit{Bangladesh Journal of Fisheries}, 35(1), 1–12. \url{https://doi.org/10.52168/bjf.2023.35.01}
\end{enumerate}

\textbf{2022 (5 publications)}
\begin{enumerate}
  \item Rouf, M. A., Islam, M. J., Roknuzzaman, M., Siddiqui, M. N., \& \textbf{Golder, M. R.} (2022). Stratification pattern of dissolved oxygen and associated water parameters in the Pasur–Rupsha estuary of Bangladesh. \textit{Heliyon}, e10935. \url{https://doi.org/10.1016/j.heliyon.2022.e10935}

  \item Bir, J., \textbf{Golder, M. R.}, \& Islam, S. S. (2022). Review on invasive alien species (IAS): Challenge and consequence to the aquatic ecosystem services. \textit{Marine Science and Technology Bulletin}, 11(3), 288–298. \url{https://doi.org/10.33714/masteb.1091625}

  \item \textbf{Golder, M. R.}, Shammi, A. T., \& Rouf, M. A. (2022). Enhanced awareness to coastal ecology: Protecting endangered species of the Bay of Bengal. \textit{Oceanography \& Fisheries Open Access Journal}, 15(1), 555904. \url{https://doi.org/10.19080/OFOAJ.2022.15.555904}

  \item Masud-Ul-Alam, M., Khan, M. A. I., Barrett, B. S., Rivera-Calle, S., \textbf{Golder, M. R.}, \& Rouf, M. A. (2022). Spatial variation of the winter thermal inversion in the northern Bay of Bengal. \textit{Regional Studies in Marine Science}, 53(6), 102417. \url{https://doi.org/10.1016/j.rsma.2022.102417}

  \item Shuva, M. S. H., \textbf{Golder, M. R.}, Rouf, M. A., Uddin, M. M., \& Bir, J. (2022). Daytime and nighttime sea surface temperature (SST) along with diurnal variability (D-SST) in the northern Bay of Bengal: A remote sensing approach. \textit{Thalassas: An International Journal of Marine Science}, 38, 929–939. \url{https://doi.org/10.1007/s41208-022-00406-8}
\end{enumerate}

\textbf{2022 continued}
\begin{enumerate}\setcounter{enumi}{5}
  \item Islam, S. S., \textbf{Golder, M. R.}, Bir, J., Mistry, S. K., Siddique, M. N. A., Ali, M. R., Sabbir, W., Azad, M. A. K., \& Huq, K. A. (2022). Growth comparison of mono and mixed-sex giant freshwater prawn (\textit{Macrobrachium rosenbergii}, De Man) and stocking optimization of male monosex farming. \textit{Khulna University Studies}, 33(2), 275–282. \url{https://doi.org/10.53808/KUS.2022.19.01.2114-ls}
\end{enumerate}

\textbf{2021 (7 publications)}
\begin{enumerate}
  \item \textbf{Golder, M. R.}, Rouf, M. A., \& Jamal, M. (2021). Variability of physical and biogeochemical parameters in the northern Bay of Bengal. \textit{Regional Studies in Marine Science}, 48, 102027. \url{https://doi.org/10.1016/j.rsma.2021.102027}

  \item \textbf{Golder, M. R.}, Zilany, M. A. K., Bir, J., Huq, K. A., Kabir, K. A., Mamun-Ur-Rashid, M., \& Islam, S. S. (2021). Effect of different experimental diets and salinity on growth performance of \textit{Penaeus monodon}. \textit{Bangladesh Journal of Fisheries}, 33(2), 275–282. \url{https://doi.org/10.52168/bjf.2021.33.30}

  \item Islam, M. R., Ali, M. Z., Hasanuzzaman, A. F. M., \& \textbf{Golder, M. R.} (2021). Effect of skip-feeding and re-feeding regimes on compensatory growth, body composition, feed utilization, feeding cost, and survival of monosex tilapia (\textit{Oreochromis niloticus}) in pond aquaculture system. \textit{Bangladesh Journal of Fisheries}, 33(2), 225–233. \url{https://doi.org/10.52168/bjf.2021.33.25}

  \item Bir, J., \textbf{Golder, M. R.}, Biswas, P., Tasrin, S., \& Chowdhury, S. Z. (2021). Plastic pollution in an estuary: A preliminary study on the Rupsha River in southwestern Bangladesh. \textit{Bangladesh Journal of Fisheries}, 33(2), 295–304. \url{https://doi.org/10.52168/bjf.2021.33.32}

  \item Rouf, M. A., \textbf{Golder, M. R.}, \& Afroje, Z. (2021). Satellite-based observation of particulate organic carbon in the northern Bay of Bengal. \textit{Environmental Advances}, 5, 100124. \url{https://doi.org/10.1016/j.envadv.2021.100124}

  \item Rouf, M. A., \textbf{Golder, M. R.}, Kana, N. A., Debnath, S., Rahman, M. M., Mathew, R. T., \& Alrashada, Y. N. (2021). Production, trading and potential utilization of fish scale in Khulna, Bangladesh. \textit{Advances in Animal and Veterinary Sciences}, 9(12), 2194–2200. \url{https://doi.org/10.17582/journal.aavs/2021/9.12.2194.2200}

  \item \textbf{Golder, M. R.}, Shuva, M. S. H., Rouf, M. A., Uddin, M. M., Bristy, S. K., \& Bir, J. (2021). Chlorophyll-a, SST and particulate organic carbon in response to Cyclone Amphan in the Bay of Bengal. \textit{Journal of Earth System Science}, 130, 157. \url{https://doi.org/10.1007/s12040-021-01668-1}
\end{enumerate}

\textbf{2021 continued}
\begin{enumerate}\setcounter{enumi}{7}
  \item \textbf{Golder, M. R.}, \& Rouf, M. A. (2021). Application of ocean color remote sensing for managing the Bay of Bengal. \textit{Bangladesh Maritime Journal}, 5(1), 25–38.
\end{enumerate}

\textbf{2020 (4 publications)}
\begin{enumerate}
  \item Bir, J., \textbf{Golder, M. R.}, Biswas, S. K., Islam, S. S., Kumar, R., \& Huq, K. A. (2020). Application of probiotics and prebiotics for promoting growth of tiger shrimp (\textit{Penaeus monodon}): An approach to eco-friendly shrimp aquaculture. \textit{International Journal of Agricultural Research, Innovation and Technology}, 10(2), 15–20. \url{https://doi.org/10.3329/ijarit.v10i2.51571}

  \item Bir, J., Das, R., \textbf{Golder, M. R.}, Islam, S. S., Paul, P., E-Mahfuj, S., Ghosh, A. K., \& Huq, K. A. (2020). Growth performance and proximate composition of grey mullet \textit{Liza parsia} at different salinities. \textit{Bangladesh Journal of Fisheries}, 32(2), 287–297. \url{https://doi.org/10.52168/bjf.2020.32.32}

  \item Bir, J., \textbf{Golder, M. R.}, \& Khalil, S. M. I. (2020). Adaptation in extreme underwater vent ecosystems: A case study on Pompeii worm (\textit{Alvinella pompejana}). \textit{International Journal of Fauna and Biological Studies}, 7(3), 25–32.

  \item Bir, J., \textbf{Golder, M. R.}, Zobayer, F. Al, Das, K. K., Zaman, S., Das, L. M., \& Paul, P. C. (2020). A review on blue economy in Bangladesh: Prospects and challenges. \textit{International Journal of Agricultural Research, Innovation and Technology}, 7(4), 21–29. \url{https://doi.org/10.5281/zenodo.4270719}
\end{enumerate}

\subsection{Press / Media}
\begin{enumerate}
  \item \textbf{Golder, M. R.} (2024). World Ocean Day: How MSP can help us navigate the blue frontier. \textit{The Daily Star}, 8 June 2024. Opinion piece. \url{https://www.thedailystar.net/opinion/views/news/world-ocean-day-how-msp-can-help-us-navigate-the-blue-frontier-3629451}
\end{enumerate}

\subsection{Peer Review Service}
\cvitem{Journals}{Remote Sensing in Earth Systems Sciences (Publons)}
\cvitem{}{Ocean Science Journal (2024)}
\cvitem{}{Geocarto International (Publons)}

\section{Technical Skills}

\subsection{Programming \& Data Science}
\cvitem{Languages}{Python (Advanced), R (Advanced), JavaScript (Intermediate)}
\cvitem{Libraries}{Pandas, NumPy, Xarray, Matplotlib, Seaborn, Plotly, ggplot2, Tidyverse}
\cvitem{Proficiency}{5+ years experience in data analysis and visualization}

\subsection{GIS \& Mapping}
\cvitem{Desktop GIS}{ArcGIS Pro (Expert), QGIS (Advanced), ENVI (Intermediate)}
\cvitem{Web GIS}{Google Earth Engine (Advanced), Leaflet, ArcGIS Online}
\cvitem{Proficiency}{7+ years experience in geospatial analysis}

\subsection{Remote Sensing}
\cvitem{Ocean Color}{SeaDAS (Expert), SNAP (Advanced), ACOLITE (Intermediate)}
\cvitem{Satellites}{MODIS, VIIRS, Sentinel-2/3, Landsat}
\cvitem{Processing}{Atmospheric correction, bio-optical algorithms, data validation}
\cvitem{Proficiency}{5+ years experience in satellite oceanography}

\subsection{Numerical Modeling}
\cvitem{Models}{TUFLOW FV (Advanced), AED - Aquatic EcoDynamics (Intermediate), Bluelink ReANalysis}
\cvitem{Skills}{Model validation, sensitivity analysis, parameter optimization}

\subsection{Database \& Tools}
\cvitem{Database}{SQL, PostgreSQL, SQLite}
\cvitem{Version Control}{Git, GitHub (5+ years)}
\cvitem{Office}{Microsoft Office Suite, LaTeX}
\cvitem{Web Dev}{HTML5, CSS3, Bootstrap, Jekyll}

\subsection{Domain Expertise}
\cvitem{Ocean Science}{Bio-optical modeling, Phytoplankton dynamics, Ocean optics}
\cvitem{GIS Analysis}{Spatial statistics, Geostatistics, Time-series analysis}
\cvitem{Data Quality}{Quality control protocols, Uncertainty quantification, Data validation}

\section{Training \& Professional Development}

\subsection{Conferences \& Workshops (8 Events, 5 Countries)}
\cvitem{Feb 2026}{Ocean Sciences Meeting (OSM2026)\\
  Glasgow, Scotland\\
  Activity: Poster Presentation\\
  Title: "Physical drivers of long-term chlorophyll-a variability in the Southern Ocean"}

\cvitem{Nov 2025}{4th Annual ARC SRI ACEAS Research Forum\\
  Hobart, Australia\\
  Activity: Conference Attendance}

\cvitem{Jun 2025}{Advancing Regional Marine Ecosystem Modelling: From Southern Ocean to Global Perspectives (FishMIP Workshop)\\
  Hobart, Australia\\
  Duration: 24–27 June 2025\\
  Activity: Workshop Participation}

\cvitem{Nov 2024}{IOCCG Summer Lecture Series 2024\\
  Hyderabad, India\\
  Duration: 4–16 November 2024 (2 weeks)\\
  Activity: Intensive Training on Ocean Color Remote Sensing\\
  Organizer: International Ocean Colour Coordinating Group (IOCCG)}

\cvitem{Nov 2023}{3rd Annual ARC SRI Australian Centre for Excellence in Antarctic Science Research Forum\\
  Hobart, Australia\\
  Duration: 7–9 November 2023\\
  Activity: Conference Attendance}

\cvitem{Aug 2023}{The Southern Ocean Observing System (SOOS) Symposium\\
  Hobart, Australia\\
  Duration: 14–18 August 2023\\
  Activity: Conference Attendance}

\cvitem{May 2021}{NOAA CoastWatch Satellite Course\\
  Online (CoastWatch, NOAA – National Oceanic and Atmospheric Administration)\\
  Activity: Online Training Course on Satellite Oceanography}

\cvitem{May--Jun 2020}{Using the Copernicus Marine Data Stream for Ocean Applications\\
  Online (EUMETSAT – European Organisation for the Exploitation of Meteorological Satellites; CMEMS – Copernicus Marine Environment Monitoring Service)\\
  Activity: Training on Copernicus Marine Data Products}

\cvitem{Nov 2019}{Ocean Colour Remote Sensing – Data, Processing and Applications\\
  Hyderabad, Telangana, India\\
  Duration: 25–29 November 2019\\
  Organizer: Indian National Centre for Ocean Information Services (INCOIS)\\
  Activity: Training on SeaDAS, Ocean color algorithms, and satellite data processing}

\subsection{Online Courses (36 Weeks, Esri Platform)}
\cvitem{2024}{Make an Impact with Modern Geo Apps (4 weeks)\\
  Provider: Esri\\
  Completed: 20 October 2024\\
  Topics: ArcGIS Instant Apps, ArcGIS Dashboards, ArcGIS StoryMaps, ArcGIS Experience Builder}

\cvitem{2024}{Cartography (6 weeks)\\
  Provider: Esri\\
  Completed: 16 May 2024\\
  Topics: Map design, Geospatial visualization techniques}

\cvitem{2024}{Going Places with Spatial Analysis (6 weeks)\\
  Provider: Esri\\
  Completed: 4 March 2024\\
  Topics: Spatial analysis tools, Problem-solving with GIS}

\cvitem{2023}{GIS for Climate Action (6 weeks)\\
  Provider: Esri\\
  Completed: 2 December 2023\\
  Topics: Climate change mitigation and adaptation using GIS}

\cvitem{2023}{Imagery in Action (6 weeks)\\
  Provider: Esri\\
  Completed: 28 October 2023\\
  Topics: Satellite and aerial imagery for GIS applications}

\cvitem{2023}{Spatial Data Science: The New Frontier in Analytics (6 weeks)\\
  Provider: Esri\\
  Completed: 7 October 2023\\
  Topics: Advanced spatial data science techniques}

\section{Ocean Cruises \& Field Work}
\cvitem{Dec 2019}{Expedition to the edge of Bangladesh Maritime Boundary in the Bay of Bengal\\
  Duration: 5–7 December 2019\\
  Vessel: BNS Turag (Bangladesh Navy Ship)\\
  Organized by: Bangladesh Navy and Khulna University\\
  Activities: Water sampling, CTD profiling, oceanographic data collection}

\cvitem{Jan 2017}{Expedition to the Swatch of No Ground in the Bay of Bengal\\
  Duration: 23–25 January 2017\\
  Vessel: BNS Surma (Bangladesh Navy Ship)\\
  Organized by: Bangladesh Navy and Khulna University\\
  Activities: Deep sea sampling, marine biodiversity survey, oceanographic observations}

\section{Professional Affiliations \& Memberships}

\cvitem{2024--Present}{Australian Marine Sciences Association (AMSA)\\
  Member}

\cvitem{2024--Present}{The Geological Society of America (GSA)\\
  Member}

\cvitem{2023--Present}{American Geophysical Union (AGU)\\
  Member ID: 1592744}

\cvitem{2022--Present}{ARC Australian Centre for Excellence in Antarctic Science (ACEAS)\\
  Member}

\cvitem{2022--Present}{The Oceanography Society (TOS)\\
  Member ID: 18077}

\cvitem{2022--Present}{Western Australian Spatial Science Student Association (WASSSA)\\
  Member}

\cvitem{2020--Present}{Krishibid Institution Bangladesh (KIB)\\
  Member ID: 06-19-31743}

\cvitem{2018--Present}{Khulna University FMRT Alumni Association\\
  Member}

\section{Research Interests \& Expertise}
\begin{itemize}
  \item Ocean Color Remote Sensing and Satellite Oceanography
  \item Ocean Optics and Bio-optical Modeling
  \item Biological Oceanography and Phytoplankton Dynamics
  \item Physical Oceanography in Polar Regions (Southern Ocean, Antarctic)
  \item Geographic Information Systems (GIS) and Spatial Analysis
  \item Environmental Science and Coastal Zone Management
  \item Marine Ecosystem Modeling and Numerical Simulations
  \item Water Quality Assessment and Monitoring
  \item Satellite Data Analysis and Algorithm Development
  \item Numerical Model Validation and Calibration
  \item Climate Change Impacts on Marine Ecosystems
  \item Biogeochemical Cycles in Aquatic Systems
  \item Machine Learning and Deep Learning for Ocean Sciences
  \item Time-series Analysis of Environmental Data
  \item Data Visualization and Scientific Communication
\end{itemize}

\section{International Collaborations}
Collaborative research with institutions across 9 countries/territories:
\begin{itemize}
  \item \textbf{Australia:} University of Tasmania, Institute for Marine and Antarctic Studies (IMAS), Curtin University, CSIRO, UWA
  \item \textbf{Bangladesh:} Bangabandhu Sheikh Mujibur Rahman Maritime University, Khulna University, University of Chittagong, WorldFish Bangladesh
  \item \textbf{USA:} Florida State University, Skidaway Institute of Oceanography, United States Air Force Academy, New Mexico State University
  \item \textbf{France:} Université de Montpellier, Université de Brest (LOPS)
  \item \textbf{Belgium:} KULeuven BIOTEC
  \item \textbf{Spain:} University of the Basque Country
  \item \textbf{China:} Tianjin University of Science \& Technology
  \item \textbf{Saudi Arabia:} King Faisal University
  \item \textbf{Thailand:} Asian Disaster Preparedness Center
\end{itemize}

\section{Languages}
\cvitemwithcomment{Bangla}{Native speaker}{}
\cvitemwithcomment{English}{Fluent (Academic \& Professional)}{}

\section{Key Recommendations}

\cvitem{Alicia Sutton}{Research Manager at WAMSI • January 23, 2026 (Mentor)\\
  \textit{"Rony would make a valued member of any team. He approaches tasks with skill, professionalism and efficiency and is always willing to tackle challenges. Rony is highly motivated to find solutions, and goes above and beyond to address needs and value add. I've enjoyed collaborating with Rony and hope to continue to in the future."}}

\cvitem{Dr. Todd Robinson}{Associate Professor at Curtin University • August 26, 2025\\
  Unit Coordinator: SPAT2008, SPAT5008, SPAT5009\\
  \textit{"Every now and then you meet a PhD student that is just that little bit extra. You know the type. They are almost always online and very responsive. They're the right mix of talent and drive. They are organised and capable. You can trust them with heavy teaching materials and know they will deliver. These guys really earn their doctorate. Rony is one of these guys."}}

\cvitem{Dr. Muhammad Abdur Rouf}{Professor and Head, Fisheries and Marine Resource Technology, Khulna University\\
  \textit{"Mr. Golder is very sound in team work and has good leading capacity. He is highly skilled in running various data management software's like R studio, ODV, ArcGIS, SeaDAS, SPSS, Python."}}

\section{References}

\subsection{Prof. David Antoine}
Head: Remote Sensing \& Satellite Research Group\\
Curtin University, School of Earth \& Planetary Sciences\\
GPO Box U1987, Perth, WA 6845, Australia\\
Associate Editor: Limnology \& Oceanography\\
\textbf{Email:} \href{mailto:david.antoine@curtin.edu.au}{david.antoine@curtin.edu.au}

\subsection{Prof. Matt Hipsey}
Professor\\
UWA School of Agriculture and Environment\\
UWA Oceans Institute, Centre for Water and Spatial Science\\
The University of Western Australia, Perth, WA 6009\\
\textbf{Email:} \href{mailto:matt.hipsey@uwa.edu.au}{matt.hipsey@uwa.edu.au}

\subsection{Dr. Alicia Sutton}
Research Manager\\
The Western Australian Marine Science Institution (WAMSI)\\
President, Australian Marine Sciences Association (AMSA)\\
35 Stirling Highway, Crawley WA 6009\\
\textbf{Email:} \href{mailto:alicia.sutton@wamsi.org.au}{alicia.sutton@wamsi.org.au}

\vspace{0.5cm}
\textit{Complete publication list and research profiles available at:}
\begin{itemize}
  \item Google Scholar: \url{https://scholar.google.com/citations?user=yXg7upIAAAAJ}
  \item ORCID: \url{https://orcid.org/0000-0002-6863-139X}
  \item ResearchGate: \url{https://www.researchgate.net/profile/Md-Golder}
  \item Personal Website: \url{https://ronygolderku.netlify.app}
\end{itemize}

\vspace{0.3cm}
\hrule
\vspace{0.2cm}
\centering
\textit{Last updated: July 2026}

\end{document}
